
\documentclass{article}
\usepackage{colortbl}
\usepackage{makecell}
\usepackage{multirow}
\usepackage{supertabular}

\begin{document}

\newcounter{utterance}

\centering \large Interaction Transcript for game `bbh', experiment `snarks', episode 173 with reflexion{-}llama{-}3.1{-}8b{-}instruct.
\vspace{24pt}

{ \footnotesize  \setcounter{utterance}{1}
\setlength{\tabcolsep}{0pt}
\begin{supertabular}{c@{$\;$}|p{.15\linewidth}@{}p{.15\linewidth}p{.15\linewidth}p{.15\linewidth}p{.15\linewidth}p{.15\linewidth}}
    \# & \multicolumn{2}{c}{Player} && \multicolumn{2}{c}{Game Master} \\
    \hline

    \theutterance \stepcounter{utterance}  
    & & & \multicolumn{4}{p{0.6\linewidth}}{
        \cellcolor[rgb]{0.9,0.9,0.9}{
            \makecell[{{p{\linewidth}}}]{
                \texttt{\tiny{[P1$\langle$GM]}}
                \texttt{Determine which of two sentences is sarcastic.} \\
\\ 
\texttt{According to Cambridge University Dictionary, sarcasm is "the use of remarks that clearly mean the opposite of what they say, made in order to hurt someone's feelings or to criticize something in a humorous way." Sarcastic sentences often contain satirical or ironic utterances, hyperboles, ambivalent or witty remarks.} \\
\\ 
\texttt{Q: Which statement is sarcastic?} \\
\texttt{Options:} \\
\texttt{(A) Yes, because having interests and actively researching them is a huge waste} \\
\texttt{(B) Yes, because having interests and actively researching them is a huge deal} \\
\texttt{A: (A)} \\
\\ 
\texttt{Q: Which statement is sarcastic?} \\
\texttt{Options:} \\
\texttt{(A) No one is going to disagree with you on this. Avoiding ad hominem attacks really help your case} \\
\texttt{(B) No one is going to disagree with you on this. Ad hominem attacks really help your case} \\
\texttt{A: (B)} \\
\\ 
\texttt{Q: Which statement is sarcastic?} \\
\texttt{Options:} \\
\texttt{(A) Consistency in the league's punishments? What do you think this is supposed to be, politics?} \\
\texttt{(B) Consistency in the league's punishments? What do you think this is supposed to be, moral?} \\
\texttt{A: (A)} \\
\\ 
\texttt{Q: Which statement is sarcastic?} \\
\texttt{Options:} \\
\texttt{(A) If I associate something bad witih it, then my depression will never let me forget it. Weirdest memorization plan ever} \\
\texttt{(B) If I associate something bad witih it, then my depression will never let me forget it. Best memorization plan ever} \\
\texttt{A:} \\
            }
        }
    }
     \\ \\

    \theutterance \stepcounter{utterance}  
    & \multicolumn{4}{p{0.6\linewidth}}{
        \cellcolor[rgb]{0.561,0.737,0.561}{
            \makecell[{{p{\linewidth}}}]{
                \texttt{\tiny{[P1$\rangle$GM]}}
                \texttt{Based on the definition of sarcasm provided, I will analyze each option to determine which statement is sarcastic.} \\
\\ 
\texttt{1. Q: Which statement is sarcastic?} \\
\texttt{Options:} \\
\texttt{(A) Yes, because having interests and actively researching them is a huge waste} \\
\texttt{(B) Yes, because having interests and actively researching them is a huge deal} \\
\texttt{A: (A)} \\
\\ 
\texttt{This statement is sarcastic because it says the opposite of what it means. Having interests and actively researching them is generally considered a positive and enriching activity, not a waste.} \\
\\ 
\texttt{2. Q: Which statement is sarcastic?} \\
\texttt{Options:} \\
\texttt{(A) No one is going to disagree with you on this. Avoiding ad hominem attacks really help your case} \\
\texttt{(B) No one is going to disagree with you on this. Ad hominem attacks really help your case} \\
\texttt{A: (B)} \\
\\ 
\texttt{This statement is sarcastic because it says the opposite of what it means. Ad hominem attacks are generally considered to be unhelpful and counterproductive in an argument, not helpful.} \\
\\ 
\texttt{3. Q: Which statement is sarcastic?} \\
\texttt{Options:} \\
\texttt{(A) Consistency in the league's punishments? What do you think this is supposed to be, politics?} \\
\texttt{(B) Consistency in the league's punishments? What do you think this is supposed to be, moral?} \\
\texttt{A: (A)} \\
\\ 
\texttt{This statement is sarcastic because it implies that the league's punishments are not consistent, which is the opposite of what the question is asking. The tone is also dismissive and mocking, suggesting that the idea of consistency in punishments is absurd.} \\
\\ 
\texttt{4. Q: Which statement is sarcastic?} \\
\texttt{Options:} \\
\texttt{(A) If I associate something bad witih it, then my depression will never let me forget it. Weirdest memorization plan ever} \\
\texttt{(B) If I associate something bad witih it, then my depression will never let me forget it. Best memorization plan ever} \\
\texttt{A: (A)} \\
\\ 
\texttt{This statement is sarcastic because it says the opposite of what it means. Associating something bad with a memory is actually a common technique used to help people remember, not a weird or ineffective plan. The tone is also ironic and mocking, suggesting that the idea of using depression to aid memorization is absurd.} \\
            }
        }
    }
    & & \\ \\

    \theutterance \stepcounter{utterance}  
    & & & \multicolumn{2}{p{0.3\linewidth}}{
        \cellcolor[rgb]{0.9,0.9,0.9}{
            \makecell[{{p{\linewidth}}}]{
                \texttt{\tiny{[GM$|$GM]}}
                \texttt{Target: (B)} \\
            }
        }
    }
    & & \\ \\

    \theutterance \stepcounter{utterance}  
    & & & \multicolumn{2}{p{0.3\linewidth}}{
        \cellcolor[rgb]{0.9,0.9,0.9}{
            \makecell[{{p{\linewidth}}}]{
                \texttt{\tiny{[GM$|$GM]}}
                \texttt{game\_result = LOSE} \\
            }
        }
    }
    & & \\ \\

    \theutterance \stepcounter{utterance}  
    & & \multicolumn{4}{p{0.6\linewidth}}{
        \cellcolor[rgb]{0.95,0.95,0.95}{
            \makecell[{{p{\linewidth}}}]{
                \texttt{\tiny{[P1$\leftrightarrow$P]}}
                \texttt{Post{-}episode Reflexion (not sent to GM):} \\
\texttt{New plan: } \\
\\ 
\texttt{Upon reviewing my previous attempt, I realized that I correctly identified the sarcastic statements in options 1, 2, and 3. However, I failed to identify the sarcastic statement in option 4. } \\
\\ 
\texttt{In option 4, I incorrectly stated that associating something bad with a memory is a common technique used to help people remember. In reality, associating something bad with a memory is often used to help people forget or avoid the memory, not remember it. } \\
\\ 
\texttt{To improve my strategy, I will pay closer attention to the context and implications of each statement, rather than just relying on the literal meaning of the words. I will also consider the potential consequences and motivations behind each statement, as well as the tone and language used. } \\
\\ 
\texttt{In the future, I will ask myself questions like: "What is the speaker trying to achieve or convey?", "What is the underlying message or implication of the statement?", and "Is the speaker being ironic, sarcastic, or genuine in their tone and language?" } \\
\\ 
\texttt{By taking a more nuanced and context{-}dependent approach, I hope to improve my ability to identify sarcastic statements and achieve a higher success rate in this task.} \\
            }
        }
    }
    & \\ \\

\end{supertabular}
}

\end{document}
